%% ma: L10-B14-0084
%% lop: 10
%% chuong: 5
%% bai: 14
%% dang: 10.14.9
%% loai: DS
%% muc: [NB, TH, TH, TH]
%% dap_an: "ĐSĐS"
%% nguon: "chuyen-de-cac-so-dac-trung-cua-mau-so-lieu-khong-ghep-nhom-toan-10-chuong-trinh-moi.pdf (GV. Phan Nhật Linh) (P05), Bài 03 Dạng 2 Phần II Câu 1"
%% kien_thuc: "Lớp 10 Bài 13 (số trung bình), Bài 14 (phương sai, độ lệch chuẩn)"
%% trang_thai: chinh-thuc
%% ghi_chu: "Tự duyệt theo yêu cầu giáo viên 22/10/2026."
\begin{ex}
Hai xạ thủ A và B, mỗi xạ thủ bắn $10$ phát đạn. Kết quả được thể hiện trong bảng sau:
\begin{center}
\begin{tabular}{|c|c|c|c|c|c|c|c|c|c|c|}\hline
A & $7$ & $9$ & $6$ & $9$ & $8$ & $6$ & $8$ & $7$ & $10$ & $8$\\\hline
B & $8$ & $7$ & $8$ & $9$ & $6$ & $7$ & $7$ & $9$ & $9$ & $8$\\\hline
\end{tabular}
\end{center}
Xét tính đúng sai của các khẳng định sau:
\choiceTF{\True Điểm thấp nhất của xạ thủ A là $6$}{Điểm trung bình của xạ thủ A cao hơn điểm trung bình của xạ thủ B}{\True Độ lệch chuẩn bảng điểm của xạ thủ A lớn hơn độ lệch chuẩn bảng điểm của xạ thủ B}{Xạ thủ A bắn đều hơn xạ thủ B}
\loigiai{a) Đúng.

b) Sai: $\overline x_A=\overline x_B=7{,}8$.

c) Đúng: $s_A^2=1{,}56$, $s_A\approx1{,}249$; $s_B^2=0{,}96$, $s_B\approx0{,}980$.

d) Sai: $s_A>s_B$ nên xạ thủ B bắn đều hơn.}
\end{ex}
